%% kapa: exam=1 round=28 year=2017 period=1 form=A subject=civil_law subject_ko=민법 seq=26 number=26
%% source: Q-Net 감정평가사 기출문제 (cst00309) artlSeq=5199603 file=extracted/1008/5199603/2017년 제28회 감정평가사 1차시험 문제지/1교시 A형.pdf
%% conversion: kordoc 4.24.51 -> tools/gen1.py (auto); review=proofread-sonnet; flags=none
\begin{question}{26}
점유에 관한 설명으로 옳지 않은 것은? (다툼이 있으면 판례에 따름)
\begin{choices}
  \item 토지매도인의 매도 후의 점유는 특별한 사정이 없는 한 타주점유로 된다.
  \item 타인소유의 토지를 자기소유 토지의 일부로 알고 이를 점유하게 된 자가 나중에 그러한 사정을 알게 되었다면 그 점유는 그 사정만으로 타주점유로 전환된다.
  \item 제3자가 토지를 경락받아 대금을 납부한 후에는 종래소유자의 그 토지에 대한 점유는 특별한 사정이 없는 한 타주점유가 된다.
  \item 토지점유자가 등기명의자를 상대로 매매를 원인으로 소유권이전등기를 청구하였다가 패소 확정된 경우, 그 사정만으로 타주점유로 전환되는 것은 아니다.
  \item 소유자가 점유자를 상대로 적극적으로 소유권을 주장하여 승소한 경우, 점유자의 토지에 대한 점유는 패소판결 확정 후부터는 타주점유로 전환된다.
\end{choices}
\end{question}
\answer{2}
